% =====================================================================
%  VERIFICA DI SISTEMI AUTOMATICI 3BETC — FILE DI RECUPERO R1 e R2.
%  Stessa struttura, punteggi e criteri della verifica (contenuto_verifica.tex),
%  esercizi diversi: banca Banca_Esercizi/sistemi/rec_*.
%  Per ottenere i PDF imposta come Main document uno dei file numerati
%  di questa cartella:
%     9_Testo_Recupero1   10_Testo_Recupero2                (da stampare)
%     11_Soluzioni_Recupero1  12_Soluzioni_Recupero2         (da dare agli alunni)
%     13_Correzione_Recupero                                 (griglia per il docente)
%  Struttura (dagli appunti di Alessio del 20/10/2026):
%     1) decimale -> binario   2) binario -> decimale
%     3-4) 3 ingressi, 2 porte logiche in cascata
%     5-6) diagrammi di flusso A SCELTA: 5 facile (6 pt) oppure 6 difficile (14 pt)
% =====================================================================
% valori predefiniti se compili direttamente questo file
\providecommand{\fila}{R1}         % R1 | R2
\providecommand{\versione}{testo}  % testo | soluzioni | correzione | completa

\documentclass[a4paper,11pt,addpoints]{exam}
\usepackage{stili/verifica}
\elencofile{R1 e R2}

% ---------- Dati della prova --------------------------------------------
\scuola{ITIS G. Galilei -- Livorno}
\annoscolastico{2026/27}
\tipoprova{Verifica}
\materia{Sistemi Automatici}
\argomento{Sistema binario, diagrammi di flusso, porte logiche}
\classe{3BETC}
\dataprova{23/10/2026}
\durata{}               % non indicata sul foglio
\votominimo{2}          % voto con 0 punti
\sogliasufficienza{60}  % % del punteggio per il 6
\criterigenerali{Conversioni: errore di calcolo isolato con procedimento corretto: metà dei punti.
  Porte logiche: colonna $Y$ 5 pt, $-1$ pt per ogni riga sbagliata, minimo 0 (non si va in negativo).
  Diagrammi di flusso: lo studente svolge \textbf{uno solo} tra l'esercizio 5 (più semplice, max 6 pt)
  e l'esercizio 6 (più difficile, max 14 pt). Punteggio massimo della prova: 42 pt con l'esercizio 6,
  34 pt con l'esercizio 5 (voto massimo circa 8). Se li svolge entrambi si conta quello con il
  punteggio più alto. Si valuta la logica dell'algoritmo; simboli sbagliati o frecce mancanti
  $-0{,}5$ pt (una sola volta); soluzioni diverse ma corrette prendono tutti i punti.}
\puntimassimi{42}       % 8+8+6+6 + max(6, 14): es. 5 e 6 a scelta
% tabella punti: es. 5 e 6 a scelta, nel totale conta solo uno dei due
\renewcommand{\tabellapunteggi}{%
  \begin{center}\small\renewcommand{\arraystretch}{1.5}%
  \begin{tabular}{|l|c|c|c|c|c|c|}\hline
    Esercizio & 1 & 2 & 3 & 4 & 5 \emph{oppure} 6 & Totale\\\hline
    Punti     & 8 & 8 & 6 & 6 & 6 \emph{oppure} 14 & 42\\\hline
    Ottenuti  &   &   &   &   &                   & \\\hline
  \end{tabular}\end{center}}

\begin{document}

\intestazione

\begin{questions}

% ---------------------------------------------------------------------
% RECUPERO (file R1, R2) - conversione decimale -> binario
%   fila R1: 179=10110011, 427=110101011, 302=100101110, 871=1101100111
%   fila R2: 478=111011110, 203=11001011, 941=1110101101, 349=101011101
%   4 numeri a 3 cifre (8, 9, 9, 10 bit), un solo pari; numeri diversi dalle file A, B, C
% ---------------------------------------------------------------------
\begin{esercizio}[8]
Converti in binario:
\RR{%
\par\smallskip\noindent\hspace*{1em}\begin{tabular}{@{}p{0.42\linewidth}p{0.42\linewidth}@{}}
  \textbf{a)}~$179_{10} = \ldots\ldots\ldots_2$ & \textbf{b)}~$427_{10} = \ldots\ldots\ldots_2$\\[4pt]
  \textbf{c)}~$302_{10} = \ldots\ldots\ldots_2$ & \textbf{d)}~$871_{10} = \ldots\ldots\ldots_2$
\end{tabular}}{%
\par\smallskip\noindent\hspace*{1em}\begin{tabular}{@{}p{0.42\linewidth}p{0.42\linewidth}@{}}
  \textbf{a)}~$478_{10} = \ldots\ldots\ldots_2$ & \textbf{b)}~$203_{10} = \ldots\ldots\ldots_2$\\[4pt]
  \textbf{c)}~$941_{10} = \ldots\ldots\ldots_2$ & \textbf{d)}~$349_{10} = \ldots\ldots\ldots_2$
\end{tabular}}

\begin{solution}
\RR{%
\noindent
  \begin{minipage}[t]{0.24\linewidth}
  \textbf{a)}\\[2pt]
  {\small\renewcommand{\arraystretch}{1.0}\begin{tabular}{r|l}
    179 & 1\\
    89 & 1\\
    44 & 0\\
    22 & 0\\
    11 & 1\\
    5 & 1\\
    2 & 0\\
    1 & 1\\
    0 & \\
  \end{tabular}}
  \end{minipage}\hfill
  \begin{minipage}[t]{0.24\linewidth}
  \textbf{b)}\\[2pt]
  {\small\renewcommand{\arraystretch}{1.0}\begin{tabular}{r|l}
    427 & 1\\
    213 & 1\\
    106 & 0\\
    53 & 1\\
    26 & 0\\
    13 & 1\\
    6 & 0\\
    3 & 1\\
    1 & 1\\
    0 & \\
  \end{tabular}}
  \end{minipage}\hfill
  \begin{minipage}[t]{0.24\linewidth}
  \textbf{c)}\\[2pt]
  {\small\renewcommand{\arraystretch}{1.0}\begin{tabular}{r|l}
    302 & 0\\
    151 & 1\\
    75 & 1\\
    37 & 1\\
    18 & 0\\
    9 & 1\\
    4 & 0\\
    2 & 0\\
    1 & 1\\
    0 & \\
  \end{tabular}}
  \end{minipage}\hfill
  \begin{minipage}[t]{0.24\linewidth}
  \textbf{d)}\\[2pt]
  {\small\renewcommand{\arraystretch}{1.0}\begin{tabular}{r|l}
    871 & 1\\
    435 & 1\\
    217 & 1\\
    108 & 0\\
    54 & 0\\
    27 & 1\\
    13 & 1\\
    6 & 0\\
    3 & 1\\
    1 & 1\\
    0 & \\
  \end{tabular}}
  \end{minipage}\\[6pt]
\par\smallskip\noindent\hspace*{1em}\begin{tabular}{@{}p{0.42\linewidth}p{0.42\linewidth}@{}}
  \textbf{a)}~$179_{10} = 10110011_2$ & \textbf{b)}~$427_{10} = 110101011_2$\\[4pt]
  \textbf{c)}~$302_{10} = 100101110_2$ & \textbf{d)}~$871_{10} = 1101100111_2$
\end{tabular}\par\smallskip
  Si divide per 2 fino ad arrivare a 0 e si leggono i resti dal basso verso l'alto.}{%
\noindent
  \begin{minipage}[t]{0.24\linewidth}
  \textbf{a)}\\[2pt]
  {\small\renewcommand{\arraystretch}{1.0}\begin{tabular}{r|l}
    478 & 0\\
    239 & 1\\
    119 & 1\\
    59 & 1\\
    29 & 1\\
    14 & 0\\
    7 & 1\\
    3 & 1\\
    1 & 1\\
    0 & \\
  \end{tabular}}
  \end{minipage}\hfill
  \begin{minipage}[t]{0.24\linewidth}
  \textbf{b)}\\[2pt]
  {\small\renewcommand{\arraystretch}{1.0}\begin{tabular}{r|l}
    203 & 1\\
    101 & 1\\
    50 & 0\\
    25 & 1\\
    12 & 0\\
    6 & 0\\
    3 & 1\\
    1 & 1\\
    0 & \\
  \end{tabular}}
  \end{minipage}\hfill
  \begin{minipage}[t]{0.24\linewidth}
  \textbf{c)}\\[2pt]
  {\small\renewcommand{\arraystretch}{1.0}\begin{tabular}{r|l}
    941 & 1\\
    470 & 0\\
    235 & 1\\
    117 & 1\\
    58 & 0\\
    29 & 1\\
    14 & 0\\
    7 & 1\\
    3 & 1\\
    1 & 1\\
    0 & \\
  \end{tabular}}
  \end{minipage}\hfill
  \begin{minipage}[t]{0.24\linewidth}
  \textbf{d)}\\[2pt]
  {\small\renewcommand{\arraystretch}{1.0}\begin{tabular}{r|l}
    349 & 1\\
    174 & 0\\
    87 & 1\\
    43 & 1\\
    21 & 1\\
    10 & 0\\
    5 & 1\\
    2 & 0\\
    1 & 1\\
    0 & \\
  \end{tabular}}
  \end{minipage}\\[6pt]
\par\smallskip\noindent\hspace*{1em}\begin{tabular}{@{}p{0.42\linewidth}p{0.42\linewidth}@{}}
  \textbf{a)}~$478_{10} = 111011110_2$ & \textbf{b)}~$203_{10} = 11001011_2$\\[4pt]
  \textbf{c)}~$941_{10} = 1110101101_2$ & \textbf{d)}~$349_{10} = 101011101_2$
\end{tabular}\par\smallskip
  Si divide per 2 fino ad arrivare a 0 e si leggono i resti dal basso verso l'alto.}
  \criterioRR
  {2 pt per ogni numero: 1 pt divisioni per 2 corrette, 1 pt resti letti dal basso e risultato. Risultati: a) $10110011_2$; b) $110101011_2$; c) $100101110_2$; d) $1101100111_2$.}
  {2 pt per ogni numero: 1 pt divisioni per 2 corrette, 1 pt resti letti dal basso e risultato. Risultati: a) $111011110_2$; b) $11001011_2$; c) $1110101101_2$; d) $101011101_2$.}
\end{solution}
\end{esercizio}

% ---------------------------------------------------------------------
% RECUPERO (file R1, R2) - conversione binario -> decimale
%   fila R1: 10111011=187, 101110011=371, 11100010=226, 1011110011=755
%   fila R2: 110001011=395, 10101001=169, 1100100011=803, 11100110=230
%   4 numeri (8, 8, 9, 10 bit), un solo pari, risultato < 1024
% ---------------------------------------------------------------------
\begin{esercizio}[8]
Converti in decimale:
\RR{%
\par\smallskip\noindent\hspace*{1em}\begin{tabular}{@{}p{0.42\linewidth}p{0.42\linewidth}@{}}
  \textbf{a)}~$10111011_{2} = \ldots\ldots\ldots_{10}$ & \textbf{b)}~$101110011_{2} = \ldots\ldots\ldots_{10}$\\[4pt]
  \textbf{c)}~$11100010_{2} = \ldots\ldots\ldots_{10}$ & \textbf{d)}~$1011110011_{2} = \ldots\ldots\ldots_{10}$
\end{tabular}}{%
\par\smallskip\noindent\hspace*{1em}\begin{tabular}{@{}p{0.42\linewidth}p{0.42\linewidth}@{}}
  \textbf{a)}~$110001011_{2} = \ldots\ldots\ldots_{10}$ & \textbf{b)}~$10101001_{2} = \ldots\ldots\ldots_{10}$\\[4pt]
  \textbf{c)}~$1100100011_{2} = \ldots\ldots\ldots_{10}$ & \textbf{d)}~$11100110_{2} = \ldots\ldots\ldots_{10}$
\end{tabular}}

\begin{solution}
\RR{%
  \textbf{a)}~$10111011_2 = 128 + 32 + 16 + 8 + 2 + 1 = 187_{10}$\\[4pt]
  \textbf{b)}~$101110011_2 = 256 + 64 + 32 + 16 + 2 + 1 = 371_{10}$\\[4pt]
  \textbf{c)}~$11100010_2 = 128 + 64 + 32 + 2 = 226_{10}$\\[4pt]
  \textbf{d)}~$1011110011_2 = 512 + 128 + 64 + 32 + 16 + 2 + 1 = 755_{10}$\\[4pt]
  Si sommano i pesi $2^0, 2^1, 2^2, \ldots$ (da destra) dei soli bit a 1.}{%
  \textbf{a)}~$110001011_2 = 256 + 128 + 8 + 2 + 1 = 395_{10}$\\[4pt]
  \textbf{b)}~$10101001_2 = 128 + 32 + 8 + 1 = 169_{10}$\\[4pt]
  \textbf{c)}~$1100100011_2 = 512 + 256 + 32 + 2 + 1 = 803_{10}$\\[4pt]
  \textbf{d)}~$11100110_2 = 128 + 64 + 32 + 4 + 2 = 230_{10}$\\[4pt]
  Si sommano i pesi $2^0, 2^1, 2^2, \ldots$ (da destra) dei soli bit a 1.}
  \criterioRR
  {2 pt per ogni numero: 1 pt pesi (potenze di 2) corretti, 1 pt somma e risultato. Risultati: a) $187_{10}$; b) $371_{10}$; c) $226_{10}$; d) $755_{10}$.}
  {2 pt per ogni numero: 1 pt pesi (potenze di 2) corretti, 1 pt somma e risultato. Risultati: a) $395_{10}$; b) $169_{10}$; c) $803_{10}$; d) $230_{10}$.}
\end{solution}
\end{esercizio}

% ---------------------------------------------------------------------
% RECUPERO (file R1, R2) - 3 ingressi, 2 porte AND/OR in cascata, un ingresso negato sulla prima porta
%   fila R1: X = AND(B, NOT C),  Y = OR(X, A)
%   fila R2: X = OR(A, NOT B),  Y = AND(X, C)
% Solo tabella di verita' (le espressioni logiche non sono ancora state spiegate)
% ---------------------------------------------------------------------
\begin{esercizio}[6]
Scrivi la tabella di verità del circuito in figura.

\begin{center}
\RR{%
  \begin{circuitikz}[x=1cm,y=1cm,scale=1.25,transform shape]
    \node[american and port] (g1) at (0,-0.3) {};
    \node[american or port] (g2) at (3,0.25) {};
    \draw (g1.in 1) -- ++(-0.9,0) node[left] {$B$};
    \draw (g1.in 2) -- ++(-0.9,0) node[left] {$C$};
    \draw (g2.in 1) -- (g1.in 1 |- g2.in 1) -- ++(-0.9,0) node[left] {$A$};
    \draw (g1.out) -- ++(0.45,0) |- (g2.in 2);
    \draw (g2.out) -- ++(0.7,0) node[right] {$Y$};
    \node[circle,draw,fill=white,inner sep=0pt,minimum size=1.6mm,anchor=east] at (g1.bin 2) {};
  \end{circuitikz}}{%
  \begin{circuitikz}[x=1cm,y=1cm,scale=1.25,transform shape]
    \node[american or port] (g1) at (0,0.3) {};
    \node[american and port] (g2) at (3,-0.25) {};
    \draw (g1.in 1) -- ++(-0.9,0) node[left] {$A$};
    \draw (g1.in 2) -- ++(-0.9,0) node[left] {$B$};
    \draw (g2.in 2) -- (g1.in 2 |- g2.in 2) -- ++(-0.9,0) node[left] {$C$};
    \draw (g1.out) -- ++(0.45,0) |- (g2.in 1);
    \draw (g2.out) -- ++(0.7,0) node[right] {$Y$};
    \node[circle,draw,fill=white,inner sep=0pt,minimum size=1.6mm,anchor=east] at (g1.bin 2) {};
  \end{circuitikz}}
\end{center}

\begin{solution}
\RR{%
\begin{center}
  \begin{circuitikz}[x=1cm,y=1cm,scale=1.25,transform shape]
    \node[american and port] (g1) at (0,-0.3) {};
    \node[american or port] (g2) at (3,0.25) {};
    \draw (g1.in 1) -- ++(-0.9,0) node[left] {$B$};
    \draw (g1.in 2) -- ++(-0.9,0) node[left] {$C$};
    \draw (g2.in 1) -- (g1.in 1 |- g2.in 1) -- ++(-0.9,0) node[left] {$A$};
    \draw (g1.out) -- ++(0.45,0) |- (g2.in 2);
    \draw (g2.out) -- ++(0.7,0) node[right] {$Y$};
    \node[circle,draw,fill=white,inner sep=0pt,minimum size=1.6mm,anchor=east] at (g1.bin 2) {};
  \end{circuitikz}
\end{center}
  {\renewcommand{\arraystretch}{1.05}\small
  \begin{tabular}{ccc|c}
    $A$ & $B$ & $C$ & $Y$\\\hline
    0 & 0 & 0 & 0\\
    0 & 0 & 1 & 0\\
    0 & 1 & 0 & 1\\
    0 & 1 & 1 & 0\\
    1 & 0 & 0 & 1\\
    1 & 0 & 1 & 1\\
    1 & 1 & 0 & 1\\
    1 & 1 & 1 & 1\\
  \end{tabular}}}{%
\begin{center}
  \begin{circuitikz}[x=1cm,y=1cm,scale=1.25,transform shape]
    \node[american or port] (g1) at (0,0.3) {};
    \node[american and port] (g2) at (3,-0.25) {};
    \draw (g1.in 1) -- ++(-0.9,0) node[left] {$A$};
    \draw (g1.in 2) -- ++(-0.9,0) node[left] {$B$};
    \draw (g2.in 2) -- (g1.in 2 |- g2.in 2) -- ++(-0.9,0) node[left] {$C$};
    \draw (g1.out) -- ++(0.45,0) |- (g2.in 1);
    \draw (g2.out) -- ++(0.7,0) node[right] {$Y$};
    \node[circle,draw,fill=white,inner sep=0pt,minimum size=1.6mm,anchor=east] at (g1.bin 2) {};
  \end{circuitikz}
\end{center}
  {\renewcommand{\arraystretch}{1.05}\small
  \begin{tabular}{ccc|c}
    $A$ & $B$ & $C$ & $Y$\\\hline
    0 & 0 & 0 & 0\\
    0 & 0 & 1 & 1\\
    0 & 1 & 0 & 0\\
    0 & 1 & 1 & 0\\
    1 & 0 & 0 & 0\\
    1 & 0 & 1 & 1\\
    1 & 1 & 0 & 0\\
    1 & 1 & 1 & 1\\
  \end{tabular}}}
  \criterioRR
  {Le 8 combinazioni di $A$, $B$, $C$ tutte presenti: 1 pt; colonna $Y$ corretta: 5 pt ($-1$ per ogni riga sbagliata, minimo 0; attenzione a NOT $C$). Colonna $Y$ attesa da 000 a 111: 00101111 (uscita della prima porta, per controllare i passaggi: 00100010).}
  {Le 8 combinazioni di $A$, $B$, $C$ tutte presenti: 1 pt; colonna $Y$ corretta: 5 pt ($-1$ per ogni riga sbagliata, minimo 0; attenzione a NOT $B$). Colonna $Y$ attesa da 000 a 111: 01000101 (uscita della prima porta, per controllare i passaggi: 11001111).}
\end{solution}
\end{esercizio}

% ---------------------------------------------------------------------
% RECUPERO (file R1, R2) - 3 ingressi, 2 porte in cascata con NAND/NOR e un ingresso negato sulla seconda porta
%   fila R1: X = NAND(A, B),  Y = AND(X, NOT C)
%   fila R2: X = NOR(B, C),  Y = OR(X, NOT A)
% Solo tabella di verita' (le espressioni logiche non sono ancora state spiegate)
% ---------------------------------------------------------------------
\begin{esercizio}[6]
Scrivi la tabella di verità del circuito in figura.

\begin{center}
\RR{%
  \begin{circuitikz}[x=1cm,y=1cm,scale=1.25,transform shape]
    \node[american nand port] (g1) at (0,0.3) {};
    \node[american and port] (g2) at (3,-0.25) {};
    \draw (g1.in 1) -- ++(-0.9,0) node[left] {$A$};
    \draw (g1.in 2) -- ++(-0.9,0) node[left] {$B$};
    \draw (g2.in 2) -- (g1.in 2 |- g2.in 2) -- ++(-0.9,0) node[left] {$C$};
    \draw (g1.out) -- ++(0.45,0) |- (g2.in 1);
    \draw (g2.out) -- ++(0.7,0) node[right] {$Y$};
    \node[circle,draw,fill=white,inner sep=0pt,minimum size=1.6mm,anchor=east] at (g2.bin 2) {};
  \end{circuitikz}}{%
  \begin{circuitikz}[x=1cm,y=1cm,scale=1.25,transform shape]
    \node[american nor port] (g1) at (0,-0.3) {};
    \node[american or port] (g2) at (3,0.25) {};
    \draw (g1.in 1) -- ++(-0.9,0) node[left] {$B$};
    \draw (g1.in 2) -- ++(-0.9,0) node[left] {$C$};
    \draw (g2.in 1) -- (g1.in 1 |- g2.in 1) -- ++(-0.9,0) node[left] {$A$};
    \draw (g1.out) -- ++(0.45,0) |- (g2.in 2);
    \draw (g2.out) -- ++(0.7,0) node[right] {$Y$};
    \node[circle,draw,fill=white,inner sep=0pt,minimum size=1.6mm,anchor=east] at (g2.bin 1) {};
  \end{circuitikz}}
\end{center}

\begin{solution}
\RR{%
\begin{center}
  \begin{circuitikz}[x=1cm,y=1cm,scale=1.25,transform shape]
    \node[american nand port] (g1) at (0,0.3) {};
    \node[american and port] (g2) at (3,-0.25) {};
    \draw (g1.in 1) -- ++(-0.9,0) node[left] {$A$};
    \draw (g1.in 2) -- ++(-0.9,0) node[left] {$B$};
    \draw (g2.in 2) -- (g1.in 2 |- g2.in 2) -- ++(-0.9,0) node[left] {$C$};
    \draw (g1.out) -- ++(0.45,0) |- (g2.in 1);
    \draw (g2.out) -- ++(0.7,0) node[right] {$Y$};
    \node[circle,draw,fill=white,inner sep=0pt,minimum size=1.6mm,anchor=east] at (g2.bin 2) {};
  \end{circuitikz}
\end{center}
  {\renewcommand{\arraystretch}{1.05}\small
  \begin{tabular}{ccc|c}
    $A$ & $B$ & $C$ & $Y$\\\hline
    0 & 0 & 0 & 1\\
    0 & 0 & 1 & 0\\
    0 & 1 & 0 & 1\\
    0 & 1 & 1 & 0\\
    1 & 0 & 0 & 1\\
    1 & 0 & 1 & 0\\
    1 & 1 & 0 & 0\\
    1 & 1 & 1 & 0\\
  \end{tabular}}}{%
\begin{center}
  \begin{circuitikz}[x=1cm,y=1cm,scale=1.25,transform shape]
    \node[american nor port] (g1) at (0,-0.3) {};
    \node[american or port] (g2) at (3,0.25) {};
    \draw (g1.in 1) -- ++(-0.9,0) node[left] {$B$};
    \draw (g1.in 2) -- ++(-0.9,0) node[left] {$C$};
    \draw (g2.in 1) -- (g1.in 1 |- g2.in 1) -- ++(-0.9,0) node[left] {$A$};
    \draw (g1.out) -- ++(0.45,0) |- (g2.in 2);
    \draw (g2.out) -- ++(0.7,0) node[right] {$Y$};
    \node[circle,draw,fill=white,inner sep=0pt,minimum size=1.6mm,anchor=east] at (g2.bin 1) {};
  \end{circuitikz}
\end{center}
  {\renewcommand{\arraystretch}{1.05}\small
  \begin{tabular}{ccc|c}
    $A$ & $B$ & $C$ & $Y$\\\hline
    0 & 0 & 0 & 1\\
    0 & 0 & 1 & 1\\
    0 & 1 & 0 & 1\\
    0 & 1 & 1 & 1\\
    1 & 0 & 0 & 1\\
    1 & 0 & 1 & 0\\
    1 & 1 & 0 & 0\\
    1 & 1 & 1 & 0\\
  \end{tabular}}}
  \criterioRR
  {Le 8 combinazioni di $A$, $B$, $C$ tutte presenti: 1 pt; colonna $Y$ corretta: 5 pt ($-1$ per ogni riga sbagliata, minimo 0; attenzione a NOT $C$). Colonna $Y$ attesa da 000 a 111: 10101000 (uscita della prima porta, per controllare i passaggi: 11111100).}
  {Le 8 combinazioni di $A$, $B$, $C$ tutte presenti: 1 pt; colonna $Y$ corretta: 5 pt ($-1$ per ogni riga sbagliata, minimo 0; attenzione a NOT $A$). Colonna $Y$ attesa da 000 a 111: 11111000 (uscita della prima porta, per controllare i passaggi: 10001000).}
\end{solution}
\end{esercizio}

\sceltaesercizi{Diagrammi di flusso: svolgi a scelta UNO SOLO tra l'esercizio 5 e l'esercizio 6.
  Con l'esercizio 5 (più semplice) il punteggio massimo della verifica è 34 punti, con
  l'esercizio 6 è 42 punti.}
% ---------------------------------------------------------------------
% RECUPERO (file R1, R2) - diagramma di flusso con calcolo e due decisioni (con suggerimenti)
%   fila R1: velocità media di un'auto -> troppo alta / troppo bassa / regolare
%   fila R2: consumo mensile di energia -> alto / basso / normale
% Stessa struttura e difficoltà dell'esercizio delle file A, B, C
% ---------------------------------------------------------------------
\begin{esercizio}[6]
\RR{%
Su un tratto di autostrada si controlla la velocità media delle auto. Disegna il diagramma di
flusso di un algoritmo che decide quale messaggio visualizzare.\\
\textbf{Dati da leggere:} $D$ = distanza percorsa (in km); $T$ = tempo impiegato (in ore).\\
\textbf{Altre variabili:} $V$ = velocità media, cioè $V = D/T$.\\
Prima controlla se $V > 130$: in questo caso visualizza ``Velocità troppo alta''. Altrimenti
controlla se $V < 60$: in questo caso visualizza ``Velocità troppo bassa''; in caso contrario
visualizza ``Velocità regolare''.}{%
Un contatore misura l'energia elettrica consumata in casa. Disegna il diagramma di flusso di un
algoritmo che giudica il consumo del mese.\\
\textbf{Dati da leggere:} $L_1$ = lettura del contatore a inizio mese (in kWh); $L_2$ = lettura a
fine mese (in kWh).\\
\textbf{Altre variabili:} $C$ = consumo del mese, cioè $C = L_2 - L_1$.\\
Prima controlla se $C > 300$: in questo caso visualizza ``Consumo alto''. Altrimenti controlla se
$C < 100$: in questo caso visualizza ``Consumo basso''; in caso contrario visualizza ``Consumo normale''.}

\begin{solution}
\begin{center}
\RR{%
\begin{tikzpicture}[x=1cm,y=1cm]
  \node[fcterm] (ini) at (0,0) {Inizio};
  \node[fcio]   (in)  at (0,-1.1) {Leggi $D$, $T$};
  \node[fcproc] (med) at (0,-2.2) {$V \gets D / T$};
  \node[fcdec]  (d1)  at (0,-3.55) {$V > 130$};
  \node[fcio]   (o1)  at (4.6,-3.55) {Scrivi ``Velocità troppo alta''};
  \node[fcdec]  (d2)  at (0,-5.25) {$V < 60$};
  \node[fcio]   (o2)  at (4.6,-5.25) {Scrivi ``Velocità troppo bassa''};
  \node[fcio]   (o3)  at (0,-6.75) {Scrivi ``Velocità regolare''};
  \node[fcterm] (fin) at (0,-8.4) {Fine};
  \draw[fcfreccia] (ini) -- (in);
  \draw[fcfreccia] (in) -- (med);
  \draw[fcfreccia] (med) -- (d1);
  \draw[fcfreccia] (d1) -- node[fcsi,above] {Sì} (o1);
  \draw[fcfreccia] (d1) -- node[fcno,right] {No} (d2);
  \draw[fcfreccia] (d2) -- node[fcsi,above] {Sì} (o2);
  \draw[fcfreccia] (d2) -- node[fcno,right] {No} (o3);
  \draw[thick] (o1.east) -- (7.4,-3.55) |- (0,-7.65);
  \draw[thick] (o2.east) -- (7.4,-5.25);
  \draw[fcfreccia] (o3) -- (fin);
\end{tikzpicture}}{%
\begin{tikzpicture}[x=1cm,y=1cm]
  \node[fcterm] (ini) at (0,0) {Inizio};
  \node[fcio]   (in)  at (0,-1.1) {Leggi $L_1$, $L_2$};
  \node[fcproc] (med) at (0,-2.2) {$C \gets L_2 - L_1$};
  \node[fcdec]  (d1)  at (0,-3.55) {$C > 300$};
  \node[fcio]   (o1)  at (4.2,-3.55) {Scrivi ``Consumo alto''};
  \node[fcdec]  (d2)  at (0,-5.25) {$C < 100$};
  \node[fcio]   (o2)  at (4.2,-5.25) {Scrivi ``Consumo basso''};
  \node[fcio]   (o3)  at (0,-6.75) {Scrivi ``Consumo normale''};
  \node[fcterm] (fin) at (0,-8.4) {Fine};
  \draw[fcfreccia] (ini) -- (in);
  \draw[fcfreccia] (in) -- (med);
  \draw[fcfreccia] (med) -- (d1);
  \draw[fcfreccia] (d1) -- node[fcsi,above] {Sì} (o1);
  \draw[fcfreccia] (d1) -- node[fcno,right] {No} (d2);
  \draw[fcfreccia] (d2) -- node[fcsi,above] {Sì} (o2);
  \draw[fcfreccia] (d2) -- node[fcno,right] {No} (o3);
  \draw[thick] (o1.east) -- (6.6,-3.55) |- (0,-7.65);
  \draw[thick] (o2.east) -- (6.6,-5.25);
  \draw[fcfreccia] (o3) -- (fin);
\end{tikzpicture}}
\end{center}
\RR{%
  Il secondo controllo si fa solo se $V \leq 130$: quindi nel ramo ``No'' basta controllare $V < 60$.}{%
  Il secondo controllo si fa solo se $C \leq 300$: quindi nel ramo ``No'' basta controllare $C < 100$.}
  \criterioRR
  {Inizio/Fine e simboli corretti (ovale, parallelogramma, rettangolo, rombo): 1 pt; lettura di $D$,
    $T$ e calcolo di $V$: 1 pt; decisione $V>130$ con i rami Sì/No: 2 pt; decisione $V<60$ nel ramo
    No: 1 pt; tre messaggi nei rami giusti e rami che si ricongiungono prima di Fine: 1 pt.}
  {Inizio/Fine e simboli corretti (ovale, parallelogramma, rettangolo, rombo): 1 pt; lettura di $L_1$,
    $L_2$ e calcolo di $C$ (attenzione all'ordine $L_2 - L_1$): 1 pt; decisione $C>300$ con i rami
    Sì/No: 2 pt; decisione $C<100$ nel ramo No: 1 pt; tre messaggi nei rami giusti e rami che si
    ricongiungono prima di Fine: 1 pt.}
\end{solution}
\end{esercizio}

% ---------------------------------------------------------------------
% RECUPERO (file R1, R2) - diagramma di flusso con ciclo (senza suggerimenti) + tabella
%   fila R1: ascensore, persone dello stesso peso -> persone salite e peso totale
%            (condizione T+P <= M, come lo scaffale della fila C)
%   fila R2: robot a batteria, compiti con consumo fisso -> compiti svolti e carica rimasta
%            (condizione B >= E, come il magazzino della fila B)
%   Prova: R1 M=300, P=70 -> 4 iterazioni;  R2 B=100, E=22 -> 4 iterazioni
% ---------------------------------------------------------------------
\begin{esercizio}[14]
\RR{%
Su un ascensore salgono persone, una alla volta, tutte dello stesso peso. Disegna il diagramma di
flusso di un algoritmo che conta quante persone possono salire.\\
\textbf{Dati da leggere:} $M$ = peso massimo che l'ascensore può portare (in kg); $P$ = peso di una
persona (in kg).\\
\textbf{Altre variabili:} $T$ = peso già sull'ascensore; $N$ = numero di persone salite.\\
Le persone salgono finché ne può salire ancora una senza superare il peso massimo $M$.
Alla fine visualizza $N$ e $T$.\\
\textbf{Prova:} con $M = 300$ e $P = 70$ completa la tabella con i valori dopo ogni iterazione del ciclo.
\begin{center}\renewcommand{\arraystretch}{1.25}
\begin{tabular}{|c|p{3cm}|p{3cm}|}\hline
  Iterazione & $T$ & $N$\\\hline
  Inizio & &\\\hline
  1 & &\\\hline
  2 & &\\\hline
  3 & &\\\hline
  4 & &\\\hline
  5 & &\\\hline
  6 & &\\\hline
\end{tabular}\end{center}}{%
Un robot svolge compiti uno dopo l'altro e ogni compito consuma sempre la stessa parte della
batteria. Disegna il diagramma di flusso di un algoritmo che conta quanti compiti riesce a svolgere.\\
\textbf{Dati da leggere:} $B$ = carica della batteria (in \%); $E$ = carica consumata da un compito (in \%).\\
\textbf{Altre variabili:} $K$ = numero di compiti svolti.\\
Il robot svolge compiti finché la batteria ha abbastanza carica per un compito intero.
Alla fine visualizza $K$ e la carica rimasta $B$.\\
\textbf{Prova:} con $B = 100$ e $E = 22$ completa la tabella con i valori dopo ogni iterazione del ciclo.
\begin{center}\renewcommand{\arraystretch}{1.25}
\begin{tabular}{|c|p{3cm}|p{3cm}|}\hline
  Iterazione & $B$ & $K$\\\hline
  Inizio & &\\\hline
  1 & &\\\hline
  2 & &\\\hline
  3 & &\\\hline
  4 & &\\\hline
  5 & &\\\hline
  6 & &\\\hline
\end{tabular}\end{center}}

\begin{solution}
\begin{center}
\RR{%
\begin{tikzpicture}[x=1cm,y=1cm]
  \node[fcterm] (ini)  at (0,0) {Inizio};
  \node[fcio]   (in)   at (0,-1.1) {Leggi $M$, $P$};
  \node[fcinit] (ini0) at (0,-2.2) {$T \gets 0$, \ $N \gets 0$};
  \node[fcdec]  (d1)   at (0,-3.75) {$T + P \leq M$};
  \node[fcproc] (agg)  at (0,-5.3) {$T \gets T + P$};
  \node[fcproc] (inc)  at (0,-6.4) {$N \gets N + 1$};
  \node[fcio]   (out)  at (5.0,-3.75) {Scrivi $N$, $T$};
  \node[fcterm] (fin)  at (5.0,-4.95) {Fine};
  \draw[fcfreccia] (ini) -- (in);
  \draw[fcfreccia] (in) -- (ini0);
  \draw[fcfreccia] (ini0) -- (d1);
  \draw[fcfreccia] (d1) -- node[fcsi,right] {Sì} (agg);
  \draw[fcfreccia] (agg) -- (inc);
  \draw[fcfreccia] (inc.south) -- ++(0,-0.5) -| (-2.4,-2.95) -- (0,-2.95);
  \draw[fcfreccia] (d1) -- node[fcno,above] {No} (out);
  \draw[fcfreccia] (out) -- (fin);
\end{tikzpicture}}{%
\begin{tikzpicture}[x=1cm,y=1cm]
  \node[fcterm] (ini)  at (0,0) {Inizio};
  \node[fcio]   (in)   at (0,-1.1) {Leggi $B$, $E$};
  \node[fcinit] (ini0) at (0,-2.2) {$K \gets 0$};
  \node[fcdec]  (d1)   at (0,-3.75) {$B \geq E$};
  \node[fcproc] (agg)  at (0,-5.3) {$B \gets B - E$};
  \node[fcproc] (inc)  at (0,-6.4) {$K \gets K + 1$};
  \node[fcio]   (out)  at (5.0,-3.75) {Scrivi $K$, $B$};
  \node[fcterm] (fin)  at (5.0,-4.95) {Fine};
  \draw[fcfreccia] (ini) -- (in);
  \draw[fcfreccia] (in) -- (ini0);
  \draw[fcfreccia] (ini0) -- (d1);
  \draw[fcfreccia] (d1) -- node[fcsi,right] {Sì} (agg);
  \draw[fcfreccia] (agg) -- (inc);
  \draw[fcfreccia] (inc.south) -- ++(0,-0.5) -| (-2.4,-2.95) -- (0,-2.95);
  \draw[fcfreccia] (d1) -- node[fcno,above] {No} (out);
  \draw[fcfreccia] (out) -- (fin);
\end{tikzpicture}}
\end{center}
\RR{%
  Ogni iterazione del ciclo è una persona che sale: $N$ conta le iterazioni. Prima di far salire una
  persona si controlla che il peso con la persona in più ($T+P$) non superi $M$; alla fine $T \leq M$.\\[4pt]
  \textbf{Prova} ($M = 300$, $P = 70$):
\begin{center}\renewcommand{\arraystretch}{1.1}
\begin{tabular}{|c|>{\centering\arraybackslash}p{2.5cm}|>{\centering\arraybackslash}p{2.5cm}|}\hline
  Iterazione & $T$ & $N$\\\hline
  Inizio & 0 & 0\\\hline
  1 & 70 & 1\\\hline
  2 & 140 & 2\\\hline
  3 & 210 & 3\\\hline
  4 & 280 & 4\\\hline
\end{tabular}\end{center}  Dopo la 4ª iterazione $T = 280$ e $280 + 70 = 350 \leq 300$ è falso: il ciclo finisce. Si visualizza $N = 4$ e $T = 280$.}{%
  Ogni iterazione del ciclo è un compito svolto: $K$ conta le iterazioni. Quando il ciclo finisce,
  nella batteria resta meno carica di quella che serve per un compito.\\[4pt]
  \textbf{Prova} ($B = 100$, $E = 22$):
\begin{center}\renewcommand{\arraystretch}{1.1}
\begin{tabular}{|c|>{\centering\arraybackslash}p{2.5cm}|>{\centering\arraybackslash}p{2.5cm}|}\hline
  Iterazione & $B$ & $K$\\\hline
  Inizio & 100 & 0\\\hline
  1 & 78 & 1\\\hline
  2 & 56 & 2\\\hline
  3 & 34 & 3\\\hline
  4 & 12 & 4\\\hline
\end{tabular}\end{center}  Dopo la 4ª iterazione $B = 12$ e $12 \geq 22$ è falso: il ciclo finisce. Si visualizza $K = 4$ e $B = 12$.}
  \criterioRR
  {Inizio/Fine e simboli corretti (inizializzazione nell'esagono): 1 pt; lettura di $M$ e $P$: 1 pt;
    inizializzazione $T\gets 0$, $N\gets 0$: 2 pt; condizione del ciclo $T+P\leq M$ ricavata da ``ne può
    salire ancora una'', con rami Sì/No giusti: 3 pt ($T<M$ o $T\leq M$, che fanno superare $M$: 1 pt);
    aggiornamento $T\gets T+P$: 2 pt; contatore $N\gets N+1$: 1 pt; ritorno alla condizione: 1 pt; uscita
    di $N$ e $T$ dopo il ciclo: 1 pt. Ciclo con condizione in fondo corretto se gestisce il caso $P>M$:
    stessi punti. Tabella della prova: 2 pt ($-0{,}5$ per ogni riga sbagliata; se il diagramma è sbagliato, vale la tabella coerente con il diagramma disegnato). Valori attesi: $T$ = 0, 70, 140, 210, 280; $N$ = 0, 1, 2, 3, 4.}
  {Inizio/Fine e simboli corretti (inizializzazione nell'esagono): 1 pt; lettura di $B$ e $E$: 1 pt;
    inizializzazione $K\gets 0$ (e nessuna inizializzazione di $B$ dopo la lettura): 2 pt; condizione del
    ciclo $B\geq E$ ricavata da ``abbastanza carica per un compito intero'', con rami Sì/No giusti:
    3 pt ($B>E$, che perde il compito quando $B=E$: 1 pt); aggiornamento $B\gets B-E$: 2 pt; contatore
    $K\gets K+1$: 1 pt; ritorno alla condizione: 1 pt; uscita di $K$ e $B$ dopo il ciclo: 1 pt. Ciclo con
    condizione in fondo corretto se gestisce il caso $B<E$ iniziale: stessi punti. Tabella della prova: 2 pt ($-0{,}5$ per ogni riga sbagliata; se il diagramma è sbagliato, vale la tabella coerente con il diagramma disegnato). Valori attesi: $B$ = 100, 78, 56, 34, 12; $K$ = 0, 1, 2, 3, 4.}
\end{solution}
\end{esercizio}


\end{questions}

% Tabella punti + griglia voti in fondo alla 2ª pagina
\riepilogo

\end{document}
